% ==============================================================================
% CMAT-HOI: Ready-to-Publish LaTeX Related Work Section
% Formatting: IEEEtran / ACM / Springer standard
% Compatible with: RELATED_PAPERS_BIBLIOGRAPHY.bib
% ==============================================================================

\section{Related Work}
\label{sec:related_work}

Automated sports action recognition lies at the intersection of video action understanding, articulated skeleton modeling, human--object interaction (HOI), and multimodal transformer architectures. This section reviews relevant literature across these four pillars and situates the design of the proposed CMAT-HOI framework.

\subsection{Video Action Recognition in Sports Analytics}
Early methods in video action recognition predominantly decoupled spatial appearance from temporal motion using Two-Stream Convolutional Networks~\cite{simonyan2014two}, which processed static RGB frames alongside precomputed dense optical flow. Carreira and Zisserman~\cite{carreira2017quo} extended 2D filters to three-dimensional spatiotemporal convolutions through the Inflated 3D ConvNet (I3D), establishing strong performance on large-scale benchmarks like Kinetics-400 and UCF-101~\cite{soomro2012ucf101}. To avoid the computational bottleneck of dense optical flow, Feichtenhofer et~al.~\cite{feichtenhofer2019slowfast} introduced SlowFast networks, deploying dual pathways operating at disparate temporal sampling frequencies. Recently, Vision Transformers have been adapted to video volumes; Arnab et~al.~\cite{arnab2021vivit} formulated ViViT using factorized spatio-temporal attention, while Liu et~al.~\cite{liu2022video} introduced Video Swin Transformer, restricting self-attention to local 3D shifted windows.

Despite their expressive power, unimodal 3D CNNs and Video Transformers suffer from pronounced \textit{scene bias} when evaluated on sports analytics~\cite{host2022overview}. In athletic events, background environments (e.g., swimming pools, tennis courts, grass pitches) often correlate strongly with specific action labels. Consequently, models frequently rely on static environmental shortcuts rather than learning the athlete's motor mechanics. As highlighted by Shao et~al.~\cite{shao2020finegym} in the FineGym benchmark, fine-grained sports movements (such as distinguishing vault rotations or subtle racket angle changes) exhibit visually identical arena backgrounds. Unimodal appearance models lack the structural inductive bias required to distinguish between actions that share identical environmental settings.

\subsection{Skeleton-Based Action Recognition and Its Limitations}
To attain invariance against varying camera angles, lighting conditions, and cluttered backgrounds, skeleton-based methods model human movement via articulated coordinate graphs. Yan et~al.~\cite{yan2018spatial} pioneered Spatial Temporal Graph Convolutional Networks (ST-GCN), defining convolutions over anatomical joint graphs along with temporal convolutions across frame sequences. Shi et~al.~\cite{shi2019two} introduced adaptive graph convolutions (2s-AGCN) to capture non-physical correlations between distant joints, while Chen et~al.~\cite{chen2021channel} refined topologies across individual channels (CTR-GCN). Duan et~al.~\cite{duan2022pyskl} proposed PoseConv3D, rasterizing 2D keypoints into stacked 3D heatmap volumes to enhance robustness against joint detection noise.

While skeleton-based architectures provide viewpoint robustness, they possess a fundamental limitation in athletic domains: \textit{object blindness}. Fine-grained sports activities are defined not merely by anatomical limb motion, but by the physical interaction between an athlete and specialized sports apparatus (e.g., racquets, bats, balls, javelins). Because standard keypoint estimators only track human anatomical landmarks, pure skeleton networks operate in complete isolation from the external athletic gear. As a result, actions with similar gross bodily kinematics---such as a cricket shot versus a baseball pitch, or a tennis serve versus a badminton smash---produce near-identical skeletal coordinate graphs, leading to misclassification.

\subsection{Human--Object Interaction (HOI) in Static and Dynamic Contexts}
Human--Object Interaction (HOI) detection seeks to localize interacting human-object pairs and classify their semantic relation triplet $\langle \text{human, predicate, object} \rangle$. Early works adopted two-stage paradigms: Gkioxari et~al.~\cite{gkioxari2018detecting} developed InteractNet, extracting human and object bounding proposals via Faster R-CNN~\cite{ren2015faster} and estimating interaction density maps. Chao et~al.~\cite{chao2018rethinking} formalized the HICO-DET benchmark, addressing long-tailed interaction distributions. With the introduction of the Transformer architecture~\cite{vaswani2017attention}, single-stage query-based approaches became dominant: Tamura et~al.~\cite{tamura2021qpic} proposed QPIC, employing direct query-based set prediction, while Kim et~al.~\cite{kim2021hotr} formulated HOTR using parallel transformer decoders.

However, the vast majority of HOI frameworks are strictly designed for static 2D images, ignoring temporal cadence and movement momentum. Dynamic athletic actions unfold over time: a tennis stroke transitions across backswing, ball impact, and follow-through. Although recent benchmarks such as Action Genome~\cite{ji2020action} and VidHOI~\cite{chao2021vidhoi} introduced spatiotemporal scene graphs for continuous video, their dense relational graph parsing incurs heavy computational latency. CMAT-HOI addresses this by formulating a targeted, lightweight Bi-Directional Cross-Attention module ($h=4, d_k=64$), allowing human joint coordinates to actively query sports apparatus detections across temporal sequences, accurately capturing physical manipulation dynamics without combinatorial graph explosion.

\subsection{Multimodal Transformers and Gradient Balancing}
Combining complementary visual, kinematic, and contextual streams is theoretically appealing, but naive feature concatenation often causes \textit{multimodal representation collapse}. As demonstrated by Wang et~al.~\cite{wang2020makes}, jointly trained multimodal networks frequently underperform their best unimodal baselines because heterogeneous modalities generalize and overfit at vastly disparate rates under unified backpropagation updates. In sports video, this disparity is particularly acute: high-dimensional visual embeddings (e.g., ResNet-50 pool5, $2048$-d) dominate over 95\% of the feature magnitude compared to pose ($85$-d) and object ($7$-d) vectors, causing gradients to concentrate almost exclusively on the visual backbone ($>93\%$ gradient norm).

Prior works have attempted to regulate multimodal interaction via directional attention (MulT~\cite{tsai2019multimodal}), latent bottleneck tokens (MBT~\cite{nagrani2021attention}), or sigmoidal gating (GMU~\cite{arevalo2017gated}). CMAT-HOI resolves multimodal representation collapse through two architectural innovations: (1) \textit{Dimension-Aligned Modality Stems} that project visual, pose, object, and kinematic streams into an identical 256-dimensional space equipped with 1D temporal convolutions, equalizing initial gradient flow; and (2) a \textit{Gated Multimodal Fusion Unit (GMFU)} that dynamically computes frame-wise softmax gating weights ($g_t \in \mathbb{R}^4, \sum g_t^{(m)} = 1$), adaptively downweighting background visual noise when kinematic and interaction cues provide decisive evidence.
